\section{Using the explicit strong trace before squaring it}
\label{sec:trace-linear}

The doubled Pell construction already exposes its strong trace as the
polynomial $f=1+2a$. Earlier compact equations nevertheless used
$f^2(cT-R)-c$ as their auxiliary ordinate. The identity $f\equiv1\pmod c$
permits the smaller expression
\[
 V=f(cT-R)-c.
\]
It preserves both residues needed by the index argument. This changes five
ten-positive-variable equations as follows; every binary addition,
subtraction and multiplication, including fixed scalar products, is charged.
The existing register $F_0=1+a$ gives $f=a+F_0$ in one addition, so every
new circuit has exactly the same operation count as its matched parent.
\begin{center}
\begin{tabular}{@{}lrrr@{}}
\toprule
Chart and transport & Parent operations/degree & New operations & New degree\\
\midrule
$q^3$ scale, $V^2$ & $75/626$ &75&402\\
$q^3$ scale, $bV$ & $76/464$ &76&352\\
$q^2$ scale, $bV$ & $79/382$ &79&290\\
Partial $q$ scale, $bV$ & $80/338$ &80&258\\
Lucas-corrected $q$ scale, $bV$ & $87/304$ &87&232\\
\bottomrule
\end{tabular}
\end{center}
These are complete polynomial specifications on the strictly positive
domain. They do not assert a global minimum or a reduction of the number
of variables. The fixed coefficient tuple is retained unchanged, including
its largest port of $592{,}884$ bits.

\subsection{The principal degree-232 equation}
The only independent unknowns are
\[
 J,F,\alpha,i,T,w,s,n,\eta,\zeta\in\mathbb Z_{>0}.
\]
Use the four whole fixed integers $\beta,K_0,M_C,M_F$ specified in the
coefficient appendix. Define polynomial abbreviations
\begin{align*}
 h&=\beta J,&q&=h+1,& C&=q-F-\alpha,\\
 X&=wq,&Y&=qs,&k&=\eta+\zeta,&c&=kY+\eta,\\
 A&=Y(X+1)+1,&\Delta&=A^2-1,&H&=4A-5,&v&=XY^2,\\
 E&=1+v(v+1)k^2,&&&&&
\end{align*}
\begin{align*}
 R&=[h(q-F)+\alpha](q^2-1)+(M_C+qM_F)J,\\
 B&=(K_0+w)C-F,\qquad U=ic^2,\qquad a=\Delta U^2E,\\
 F_0&=1+a,\qquad G=EF_0,\qquad f=1+2a,\qquad Q=4\Delta aF_0,\\
 b&=cT-R,\qquad V=fb-c,\qquad L=bV,\qquad K=X(3c-2X)+2,\\
 D&=6q^6,\qquad
 N=6(k-R-1-q^4s)-v(R+1)(R-1)(R+3),\\
 M&=hDH,\\
 Z&=hD[H(n-LB)+K]-H[DB+hN].
\end{align*}
The single equation is
\begin{equation}
 \boxed{\ P=Q\bigl(Z^2-M^2GV^2\bigr)+M^2G-Z^2=0.\ }
 \label{eq:trace232}
\end{equation}
No intermediate register is quantified. All operations are addition,
subtraction and multiplication of integer polynomials.

\subsection{Why the trace replacement is sound}
The key identities are
\[
 f^2=1+4aF_0,\qquad
 Q=\Delta(f^2-1)=4\Delta^2U^2G,\qquad
 f-1=2\Delta i^2c^4E.
\]
At a positive zero, put $z=Z/M$. The final equation reads
$(Q-1)z^2=G(QV^2-1)$. Because $G$ is relatively prime to both other
factors, prime valuations force $G=g^2$; its relatively prime positive
factors $E$ and $F_0$ are separately squares. The adjacent-coefficient
Pell descent makes $y=|z|/g$ integral, and
\[
 QV^2-(Q-1)y^2=1,\quad
 f^2-\Delta I^2=1,\quad I=2Ug,\quad c^2\mid I.
\]
The signed folding theorem uses $V\equiv-c\pmod f$ to recover
$c=\psi_A(p)$ with odd $p>0$. The first Pell equation yields
$k=2\psi_{2XY^2+1}(n_0)$ and $n_0\ge2$, hence
$c>8XY^3\ge8q^4$. Main index three is excluded by the unchanged
quotient-strip argument: it would force $2Y-X\in\{1,2\}$ although
$q\mid X,Y$ and $q\ge32$.

One must also recheck the auxiliary index and both possible signs of $V$.
Let $c_3=4\Delta+3$. Now $c>c_3$ and $f>2c$. If the auxiliary index
were three, $V=\pm(4Q-3)$ would give one of
\[
 f(b-4\Delta f)=c-c_3,\qquad
 f(b+4\Delta f)=c+c_3.
\]
Each right side is strictly between zero and $f$. Thus the auxiliary
index is at least five, $|V|>8Q^2$, and
$|b|>(8Q^2-c)/f>8f^3-1/2>4g$.
Consequently $L=bV>0$ and $gy<L/2$. These estimates do not assume
that either $b$ or $V$ is already positive.

Writing $t=B/h$, $\kappa=K/H$, $\sigma=s/q^2$ and
\[
 \rho=\frac{k-R-1-v(R+1)(R-1)(R+3)/6}{q^6},
\]
the exact rational carrier is
$z=n+\kappa+\sigma-(1+hL)t-\rho$.
The signed transport proof first gives $B\ge0$ and $C>0$.
The only adverse cubic cases are $R=-r\le-4$ and $R=0$.
For $R=-r\le-4$, the mask gives $|C|>r/q^3$ and
$B\le-1-w|C|$. Since $b>c,r$, the new trace still gives
$V>(f-1)b$ and $L>(f-1)r^2$. The inequality
$f-1=2a>s^2$ makes $w|C|L$ dominate the entire adverse cubic term
$ws^2r^3/(6q^3)$. At $R=0$ the same estimate dominates
$ws^2/(2q^3)$. Thus both exceptional cases are covered without
assuming the signs they are meant to prove.

Only after $C>0$ is established does one use $3q<R<q^4$ and obtain
$c>2R$. The literal trace congruence gives
\[
 V\equiv-c\pmod f,\qquad V\equiv-R\pmod c.
\]
These are exactly the hypotheses of the signed auxiliary-index theorem,
so $p=R$. The retained Lucas correction then authenticates
$2n_0=R+1$, restores $q^2\mid s$, and proves that $\rho$ is integral.
The whole-product Pell projection forces $X=2^R$; the remaining carrier
gives $h\mid B$. The inherited parity, half-binomial and native proofs now
have their full hypotheses, and give a finite ZFC refutation. The complete
argument, including the all-prime corrected-index divisibility, is in
\texttt{compact/PROOF.md} and its retained pinned sources.

\subsection{An exact polynomial completeness map}
There is a useful stronger statement than a fresh canonical construction.
Put $d=2\Delta i^2c^3E$, so $f=1+cd$. From every positive zero of the
matched parent define
\[
 T_{\rm new}=fT_{\rm old}-dR.
\]
Then $b_{\rm new}=fb_{\rm old}$ and
$V_{\rm new}=f^2b_{\rm old}-c=V_{\rm old}$ as polynomial identities.
For the $bV$ transport choose
\[
 n_{\rm new}=n_{\rm old}+(f-1)L_{\rm old}B.
\]
This preserves $n-LB$, hence $Z$ and the entire final polynomial.
For the $V^2$ transport retain $n_{\rm new}=n_{\rm old}$.
The parent zero theorem gives $R<c/2$, $B\ge0$ and positive
$b_{\rm old},V_{\rm old}$, so
$T_{\rm new}>f-(f-1)/2>0$ and $n_{\rm new}\ge n_{\rm old}>0$.
This is a polynomial map from every parent positive zero; no integral
inverse on all new zeros is asserted. Combined with the soundness proof,
it establishes the desired existence equivalence.

\subsection{Attained degree and finite evidence}
Put $q_0=\beta J$, $k_0=\eta+\zeta$, and
$B_0=k_0sT-q_0^2(q_0-F)$. The entire highest homogeneous part is
\[
 [P]_{232}=-9\,2^{10}q_0^{100}w^{26}s^{58}i^{10}k_0^{32}B_0^2.
\]
It is nonzero at the actual coefficient tuple: the coefficient of
$J^{100}w^{26}s^{60}i^{10}\eta^{34}T^2$ is
$-9\,2^{10}\beta^{100}\ne0$. The degree tie between $cT$ and $R$ is
retained in $B_0$.
The supplied independent algebra checker binds every actual circuit row to
its formula, authenticates all fixed numerals, proves every whole leading
form, checks the exact parent transformation, and rejects actual gate
corruptions. These finite checks supplement the unbounded proof and its
named native/refutation calibration. They are not a newly compiled
end-to-end proof-assistant certificate.
